% !PRINTED 295
\weberpage{295}
%<<CONT
the possibility of determining the characteristic completely from
the signs of the functions $\varphi, \psi$ at the points of division of the
segments. But how we are to divide the curve $f$ into such subsegments
is taught us by the \gesperrt{Sturm theorem}. For this purpose we imagine
$x$ and $y$ represented as functions of a variable $t$, which runs along the curve $f$ through all
real values from $-\infty$ to $+\infty$
(or also only through the values of a
finite interval). Thereby $\varphi$ and $\psi$ also pass into functions of the variable $t$, and one
has only to investigate the position of their roots according to the \gesperrt{Sturm theorem}.

\weberfigureimage{weber-b1-fig14.png}{0.44}

If, for example, as in Fig. 14, the curve $f$ is a circle with the
equation
\[
(x - \alpha)^2 + (y - \beta)^2 = \varrho^2,
\]
then we can set:
\[
x - \alpha = \varrho\, \frac{1 - t^2}{1 + t^2},
\]
\[
y - \beta = \varrho\, \frac{2t}{1 + t^2},
\]
and $t$ simply runs through all values from $-\infty$ to $+\infty$, while
the point $x, y$ moves over the whole circle in the positive sense.
$\varphi$ and $\psi$, multiplied by a suitable power of
$1 + t^2$, pass into integral rational functions of $t$,
which permit the application of the \gesperrt{Sturm theorem}.

\begin{center}
\S.\,98.

\gesperrt{Gauss' first proof of the fundamental theorem of algebra.}
\end{center}

\gesperrt{Gauss} has given three different proofs of the principal theorem
of algebra, that an equation of degree $n$ has $n$ roots.
The first of these, which is found in his doctoral dissertation, and which he later carried out still further and
presented differently, rests, though in a different garb,
%CONT>>
